\chapter{Что вычисляют \nt{ГЛУТ}-нейроны}
\label{sec:glutcomputer}

\section{Правила игры}

В главе~\ref{sec:neurontypes} мы увидели, что подавляющее большинство нейронов мозга --- \nt{ГЛУТ}: они только возбуждают, их веса только положительны. Возьмём сколько угодно таких нейронов и спросим: что может вычислить такая сеть, если разрешать себе только то, что бывает в живом организме?

А в живом организме нет тактового генератора, который переключал бы все элементы одновременно. Каждый нейрон работает сам по себе, в непрерывном времени: импульсы приходят и уходят когда придётся, с разными задержками. Есть \emph{входы} --- чувствительные нейроны, которые срабатывают от света, звука или давления, --- и \emph{выходы} --- нейроны, которые управляют мышцами. Между ними --- сеть, и никто не говорит ей, когда начинать и когда заканчивать вычисление. Для электронщика это \emph{асинхронная} схема с элементами, задержки которых заранее точно не известны.

Модель нейрона возьмём самую простую из тех, что честно работают в непрерывном времени. Напряжение $V$ на мембране нейрона в покое равно нулю. Каждый пришедший импульс от нейрона $j$ скачком поднимает его на $w_j\ge0$, после чего напряжение само стекает обратно к нулю с постоянной времени $\tau$:
\begin{empheq}[box=\keybox]{equation}
\tau\,\frac{\dd V}{\dd t} \;=\; -V \;+\; \tau\sum_j w_j \sum_k \delta\bigl(t-t_j^{(k)}-d_j\bigr), \qquad w_j\ge 0 .
\label{eq:glutneuron}
\end{empheq}
Здесь $t_j^{(k)}$ --- моменты импульсов нейрона $j$, $d_j$ --- задержка на пути от него, $\delta$ --- дельта-функция, которая и означает мгновенный скачок. Когда $V$ достигает порога $\theta$, нейрон выдаёт импульс, $V$ сбрасывается в ноль, и примерно миллисекунду нейрон не может сработать снова. Типичные числа: $\tau$ --- от долей миллисекунды до двадцати миллисекунд у разных нейронов, задержки --- от долей миллисекунды до десятков миллисекунд. Это \emph{интегрирующий нейрон с утечкой}: RC-цепь с порогом. Это сознательное упрощение: у нейронов коры ветви входов сами выдают местные импульсы~\cite{Smith2013}, и один нейрон ведёт себя скорее как небольшая многослойная сеть (глава~\ref{sec:synthesis}). Для вопросов этой главы хватает и одной RC-цепи.

Главное отличие от логического вентиля --- утечка: нейрон помнит импульс не вечно, а примерно $\tau$, и складывается только то, что пришло в пределах этого окна.

\section{ИЛИ и И}

Начнём с вентилей. Если один импульс поднимает напряжение выше порога, $w\ge\theta$, нейрон срабатывает на каждый импульс любого входа. Это \emph{ИЛИ}.

С И интереснее. Пусть $w<\theta<2w$: одного импульса мало, нужны два. Но теперь вопрос «пришли ли оба» не имеет смысла, пока не сказано, \emph{в пределах какого времени}. Пусть первый импульс пришёл в момент $0$, а второй --- через время $\Delta$. К приходу второго от первого осталось $we^{-\Delta/\tau}$, и нейрон сработает, если $we^{-\Delta/\tau}+w\ge\theta$. Отсюда окно совпадения:
\begin{empheq}[box=\keybox]{equation}
\Delta \;\le\; \Delta_{\max} \;=\; \tau\,\ln\frac{w}{\theta-w} .
\label{eq:window}
\end{empheq}
Это не логическое И, а \emph{детектор совпадений}: И во времени (рис.~\ref{fig:coincidence}). При $\theta=1.5w$ окно равно $\tau\ln2\approx0.7\tau$. У нейрона коры с $\tau=10\ms$ это около семи миллисекунд. У нейронов слуховой системы, где $\tau$ меньше миллисекунды, окно сжимается до долей миллисекунды.

\begin{figure}[t]
\centering
\begin{tikzpicture}[>=Stealth,x=0.9cm,y=1.6cm]
\foreach \dx/\lab/\c [count=\i] in {0.5/{совпали}/red!70!black, 2.6/{не совпали}/blue!60!black}{
 \begin{scope}[xshift={(\i-1)*7.2cm}]
  \draw[->,gray!70!black] (0,0) -- (6.3,0) node[below left,black] {\small $t$};
  \draw[->,gray!70!black] (0,0) -- (0,1.75) node[left,black] {\small $V$};
  \draw[densely dashed,gray!60!black] (0,1.5) -- (6.0,1.5) node[right,black] {\small $\theta$};
  \draw[very thick,\c] (0,0) -- (1,0) -- (1,1)
      plot[domain=1:1+\dx,samples=30] (\x,{exp(-(\x-1)/1.5)});
  \pgfmathsetmacro{\v}{exp(-\dx/1.5)+1}
  \draw[very thick,\c] (1+\dx,{exp(-\dx/1.5)}) -- (1+\dx,{min(\v,1.5)});
  \ifdim\v pt<1.5pt
    \draw[very thick,\c] plot[domain=1+\dx:6,samples=40] (\x,{\v*exp(-(\x-1-\dx)/1.5)});
  \else
    \draw[very thick,\c] (1+\dx,1.5) -- (1+\dx,0) -- (6,0);
    \draw[->,thick,\c] (1+\dx,1.55) -- (1+\dx,1.72) node[above,font=\scriptsize] {импульс};
  \fi
  \draw[->,thick] (1,-0.35) -- (1,-0.08); \draw[->,thick] (1+\dx,-0.35) -- (1+\dx,-0.08);
  \draw[decorate,decoration={brace,amplitude=3pt,mirror}] (1,-0.42) -- (1+\dx,-0.42) node[midway,below=3pt] {\scriptsize $\Delta$};
  \node[\c] at (3.5,-0.95) {\small \lab};
 \end{scope}}
\end{tikzpicture}
\caption{Детектор совпадений, порог $\theta=1.5w$. Слева второй импульс пришёл через $\Delta<\Delta_{\max}$, сумма достигла порога, нейрон сработал. Справа $\Delta>\Delta_{\max}$: от первого импульса почти ничего не осталось, и нейрон промолчал.}
\label{fig:coincidence}
\end{figure}

Другой способ получить И --- через длительность, а не через точное время. Пока горит свет, чувствительный нейрон выдаёт импульсы часто и ровно; если порогу мало одного такого входа, а двух хватает, нейрон работает, пока активны оба. Так сеть вычисляет по \emph{уровням}, и точное время прихода импульсов неважно. Большая часть того, что ниже, будет именно так.

\section{Чего сделать нельзя}

Теперь попробуем сделать НЕ: нейрон, который работает, когда вход молчит. Не получится: без входных импульсов в~\eqref{eq:glutneuron} нет ни одного скачка, и напряжение стоит на нуле, ниже порога. А сеть из многих нейронов? Тоже нет, и это доказывается сразу для всех сетей.

Удобнее всего говорить о частотах. Пусть $r_i$ --- частота импульсов нейрона $i$, $I_i$ --- приток от входов, а $f_i$ --- его передаточная характеристика: как частота зависит от суммарного притока. У интегрирующего нейрона эта характеристика неубывающая: больше приток --- чаще импульсы. Частоты в сети, пришедшей в равновесие, удовлетворяют уравнениям
\begin{equation}
r_i \;=\; f_i\Bigl(\sum_j w_{ij}\,r_j + I_i\Bigr), \qquad w_{ij}\ge 0 .
\label{eq:ratenet}
\end{equation}

\begin{keyblock}
\begin{proposition}[Монотонность]
\label{prop:monotone}
Пусть сеть~\eqref{eq:ratenet} состоит только из \nt{ГЛУТ}-нейронов и начинает с полной тишины. Если усилить входы, то установившаяся частота ни одного нейрона не уменьшится. Всё, что вычисляет такая сеть, --- \emph{монотонные} функции входов.
\end{proposition}
\end{keyblock}

Доказательство. Начнём с тишины, $r^{(0)}=0$, и будем подставлять частоты в правую часть~\eqref{eq:ratenet}: $r^{(k+1)}=F\bigl(r^{(k)}\bigr)$. Правая часть $F$ не убывает ни по одному аргументу, потому что веса неотрицательны, а $f_i$ не убывают. Поскольку $r^{(1)}\ge0=r^{(0)}$, по индукции $r^{(k+1)}\ge r^{(k)}$: частоты только растут и приходят к наименьшему решению~\eqref{eq:ratenet}. Если входы $I$ заменить на большие $I'$, то на каждом шаге $r^{(k)}(I)\le r^{(k)}(I')$, и в пределе тоже. Настоящая сеть приходит в равновесие не шагами, а непрерывно, но для неё верно то же самое: при положительных связях неравенство между двумя прогонами, выполненное вначале, сохраняется всё время.

Для отдельных импульсов утверждение верно не буквально: лишний входной импульс может заставить нейрон сработать чуть раньше, и из-за миллисекунды невосприимчивости пропадёт его следующий импульс. Но эффект слабый и ненадёжный, вычислять на нём нельзя. Для частот монотонность строгая.

Следствия те же, что у любой схемы без инверторов. Нет НЕ. Нет исключающего ИЛИ: добавление второго входа не может погасить выход. И самое неприятное: \emph{активность нельзя остановить активностью}, можно только убрать то, что её поддерживает.

\section{Два провода: включение и выключение}

Выход из тупика подсказывает сам организм. Во многих органах чувств раздражитель передаётся двумя наборами нейронов: в зрении одни клетки отвечают на усиление света, другие --- на ослабление~\cite{Schiller1992}. Назовём их \emph{включающими} и \emph{выключающими}. Вместо одного провода $x$ сеть получает пару $(x,\bar x)$, и отсутствие, которое она не умеет вычислить, ей сообщает сам датчик.

С парой проводов НЕ бесплатно: достаточно поменять провода местами. И и ИЛИ делаются каждое двумя нейронами:
\begin{empheq}[box=\keybox]{equation}
\begin{aligned}
x\wedge y \;&\to\; \bigl(\;x\wedge y,\;\; \bar x\vee\bar y\;\bigr),\\
x\vee y \;&\to\; \bigl(\;x\vee y,\;\; \bar x\wedge\bar y\;\bigr).
\end{aligned}
\label{eq:dualrail}
\end{empheq}
Справа все операции монотонны, так что утверждение~\ref{prop:monotone} не нарушено.

Для асинхронной схемы у двухпроводного кода есть второе достоинство, даже более важное. Пара проводов бывает в трёх состояниях: «да», «нет» и, когда оба молчат, \emph{«ответа пока нет»}. Получатель узнаёт, что вычисление закончено, не по часам, а по самому сигналу: в каждой паре заработал ровно один провод. Между раздражителями датчики молчат, и сеть возвращается в тишину, готовая к следующему разу. В асинхронной электронике на этом приёме строят схемы, которые работают правильно при любых задержках элементов~\cite{Sparso2001}.

\begin{keyblock}
\begin{proposition}[Асинхронное вычисление]
\label{prop:dualrail}
Если каждый вход приходит парой проводов «включение---выключение», то сеть из \nt{ГЛУТ}-нейронов может вычислить любую логическую функцию входов. Ответ тоже приходит парами проводов, а его готовность видна по тому, что в каждой паре заработал один провод. Часы для этого не нужны. Цена --- примерно вдвое больше нейронов.
\end{proposition}
\end{keyblock}

Пример --- исключающее ИЛИ: «изменился ровно один из двух раздражителей». Прямой провод ответа --- $(x\wedge\bar y)\vee(\bar x\wedge y)$, обратный --- $(x\wedge y)\vee(\bar x\wedge\bar y)$. Это шесть нейронов, и ни один из них не нуждается в торможении (рис.~\ref{fig:dualxor}).

\begin{figure}[t]
\centering
\begin{tikzpicture}[>=Stealth,x=1cm,y=1cm,
  nrn/.style={circle,draw,very thick,fill=blue!6,minimum size=0.75cm,inner sep=0pt,font=\small},
  inp/.style={font=\small}]
\node[inp] (X)  at (0,3.3) {$x$};
\node[inp] (Xb) at (0,2.2) {$\bar x$};
\node[inp] (Y)  at (0,1.1) {$y$};
\node[inp] (Yb) at (0,0)   {$\bar y$};
\node[nrn] (A1) at (3.2,3.3) {$2$};
\node[nrn] (A2) at (3.2,2.2) {$2$};
\node[nrn] (A3) at (3.2,1.1) {$2$};
\node[nrn] (A4) at (3.2,0)   {$2$};
\node[nrn] (O)  at (7.2,2.75) {$1$};
\node[nrn] (Ob) at (7.2,0.55) {$1$};
\foreach \s/\t in {X/A1,Yb/A1,Xb/A2,Y/A2,X/A3,Y/A3,Xb/A4,Yb/A4}{\draw[->,thick,blue!60!black] (\s) -- (\t);}
\draw[->,thick,blue!60!black] (A1) -- node[pos=0.45,above,sloped,font=\scriptsize,black] {$x\wedge\bar y$} (O);
\draw[->,thick,blue!60!black] (A2) -- node[pos=0.45,below,sloped,font=\scriptsize,black] {$\bar x\wedge y$} (O);
\draw[->,thick,blue!60!black] (A3) -- node[pos=0.45,above,sloped,font=\scriptsize,black] {$x\wedge y$} (Ob);
\draw[->,thick,blue!60!black] (A4) -- node[pos=0.45,below,sloped,font=\scriptsize,black] {$\bar x\wedge\bar y$} (Ob);
\draw[->,very thick,red!70!black] (O) -- (8.6,2.75) node[right,black] {$x\oplus y$: «да»};
\draw[->,very thick,red!70!black] (Ob) -- (8.6,0.55) node[right,black] {$\overline{x\oplus y}$: «нет»};
\node[below,font=\small,gray!40!black] at (3.2,-0.55) {И: порог $2$};
\node[below,font=\small,gray!40!black] at (7.2,-0.55) {ИЛИ: порог $1$};
\node[below,font=\small,gray!40!black] at (0,-0.55) {пары проводов};
\end{tikzpicture}
\caption{Исключающее ИЛИ на двухпроводном коде, только из \nt{ГЛУТ}-нейронов. Число в кружке --- порог в единицах веса одного входа. Четыре нейрона И узнают четыре сочетания входов, два нейрона ИЛИ собирают их в пару проводов ответа.}
\label{fig:dualxor}
\end{figure}

\section{Задержки вместо часов}

Для вычислений со временем достаточно задержек в проводах и детекторов совпадений. Лучший пример --- то, как мозг определяет, откуда пришёл звук.

Звук от источника, расположенного сбоку, приходит в ближнее ухо раньше, чем в дальнее. Разница зависит от направления: от нуля, когда источник прямо впереди, до $0.22\,\text{м}/343\,\text{м/с}\approx0.6\ms$ у человека, когда он точно сбоку. Мозг различает разницу около десяти \emph{микросекунд}~\cite{KlumppEady1956,Grothe2010}, хотя сами нейроны срабатывают с точностью не лучше долей миллисекунды. Как?

Проведём от левого уха провод вдоль ряда детекторов совпадений слева направо, а от правого --- справа налево (рис.~\ref{fig:itd}). Сигнал бежит по проводу со скоростью $v$. До детектора, стоящего на расстоянии $x$ от левого края ряда длины $L$, левый сигнал идёт время $x/v$, правый --- $(L-x)/v$. Если в правое ухо звук пришёл позже на $\delta t$, два сигнала встретятся одновременно в той точке, где
\begin{empheq}[box=\keybox]{equation}
\frac{x}{v} \;=\; \frac{L-x}{v} + \delta t
\quad\Longrightarrow\quad
x \;=\; \frac{L}{2} + \frac{v\,\delta t}{2} .
\label{eq:itd}
\end{empheq}
Сработает только тот детектор, к которому оба сигнала пришли в пределах окна~\eqref{eq:window}. Номер сработавшего детектора и есть направление на источник. Разница во времени превратилась в \emph{место}.

\begin{figure}[t]
\centering
\begin{tikzpicture}[>=Stealth,
  nrn/.style={circle,draw,very thick,fill=blue!6,minimum size=0.7cm,inner sep=0pt}]
\foreach \k in {0,...,6}{\node[nrn] (D\k) at (1.4*\k,0) {\scriptsize $2$};}
\node[nrn,fill=red!15] at (D4) {\scriptsize $2$};
\draw[->,thick,blue!60!black] (-1.4,0.9) node[left,black] {\small левое ухо} -- (9.2,0.9);
\draw[->,thick,orange!85!black] (9.8,-0.9) node[right,black] {\small правое ухо} -- (-0.8,-0.9);
\foreach \k in {0,...,6}{
  \draw[->,blue!60!black] (1.4*\k,0.9) -- (D\k.north);
  \draw[->,orange!85!black] (1.4*\k,-0.9) -- (D\k.south);}
\draw[->,thick,red!70!black] (D4.north east) ++(0.1,0.05) -- ++(0.5,0.45) node[above right,font=\small,black] {сработал};
\node[font=\small,gray!40!black] at (4.2,-1.5) {задержка растёт $\longrightarrow$ \hspace{2.2cm} $\longleftarrow$ задержка растёт};
\pic[scale=1.4] at (-2.3,1.75) {speaker};
\node[font=\scriptsize,align=center] at (-2.3,2.35) {источник слева};
\end{tikzpicture}
\caption{Определение направления на звук. Сигналы от двух ушей бегут навстречу друг другу; каждый кружок --- детектор совпадений с порогом $2$. Если звук пришёл в правое ухо позже, сигналы встречаются правее середины, по формуле~\eqref{eq:itd}. Выход сети --- номер сработавшего нейрона. Здесь источник слева: в левое ухо звук пришёл раньше.}
\label{fig:itd}
\end{figure}

Здесь нет ни часов, ни торможения, ни даже двухпроводного кода: сеть не отвечает «да» или «нет», а указывает на один нейрон из многих. Такая схема из линий задержки и детекторов совпадений, по-видимому, действительно работает в слуховом стволе птиц~\cite{Carr1990}. Микросекундная точность получается не из точности каждого нейрона, а из того, что детекторов много и каждый смотрит на свою задержку. У млекопитающих ряда таких задержек не нашли: там ту же задачу, по-видимому, решают детекторы совпадений, настроенные точно рассчитанным во времени торможением от \nt{ГЛИЦ}-нейронов~\cite{Brand2002,Grothe2010}. Как торможение сужает окно совпадения, мы разберём в главе~\ref{sec:gaba} на примере \nt{ГАМК}; глицин тормозит так же, открывая у получателя хлорные каналы.

\section{Память}

Компьютеру нужна память. В такой сети память --- это активность, которая поддерживает сама себя. Возьмём группу \nt{ГЛУТ}-нейронов, сильно связанных друг с другом, и опишем её одной частотой $r$ (рис.~\ref{fig:recurrentgroup}а). В равновесии $r=f(wr+I)$, где $w$ --- сила обратной связи группы на саму себя, а $I$ --- приток извне. Решим это уравнение графически, пересекая кривую $f(wr+I)$ с прямой $r$ (рис.~\ref{fig:bistable}).

\begin{figure}[t]
\centering
% (b): tau dr/dt = -r + f(w r + I), f as in fig:bistable, w=1, I=0.6; Euler dt=0.001 tau.
\begin{tikzpicture}[>=Stealth,
  nrn/.style={circle,draw,thick,fill=blue!6,minimum size=0.5cm,inner sep=0pt}]
\begin{scope}
\node[font=\small] at (-0.5,1.55) {(а)};
\foreach \k in {1,...,6}{\node[nrn] (n\k) at ({1.4+1.0*cos(60*\k)},{1.0*sin(60*\k)}) {};}
\foreach \a/\b in {1/2,2/3,3/4,4/5,5/6,6/1,1/4,2/5,3/6}{\draw[<->,blue!60!black] (n\a) -- (n\b);}
\foreach \k in {2,3,4}{\draw[->,thick,green!45!black] ($(n\k)+(-0.95,0)$) -- (n\k);}
\node[font=\small,green!45!black] at (-0.35,-0.45) {$I$};
\node[font=\Large] at (3.1,0) {$\approx$};
\node[nrn,minimum size=0.85cm,font=\small] (R) at (4.35,-0.1) {$r$};
\draw[->,thick,blue!60!black] (R) edge[loop above,looseness=6] node[above,font=\small] {$w$} (R);
\draw[->,thick,green!45!black] (3.55,-0.95) node[left,font=\small] {$I$} -- (R);
\end{scope}
\begin{scope}[xshift=6.3cm,y=2.6cm,x=1.05cm]
\node[font=\small] at (-0.6,1.25) {(б)};
\draw[->,gray!70!black] (0,0) -- (7.3,0) node[below left,font=\small,black] {$t/\tau$};
\draw[->,gray!70!black] (0,0) -- (0,1.12) node[left,font=\small,black] {$r$};
\foreach \t in {0,2,4,6}{\draw[gray!60!black] (\t,-0.02) -- (\t,0.02); \node[below,font=\scriptsize] at (\t,0) {\t};}
\draw[densely dotted,gray!60!black] (0,0.5) -- (7,0.5) node[right,font=\scriptsize,black,align=left] {порог\\записи};
\fill[green!45!black,opacity=0.25] (1,-0.2) rectangle (2,-0.13);
\fill[gray!60!black,opacity=0.35] (1,-0.29) rectangle (1.5,-0.22);
\draw[very thick,blue!60!black] plot coordinates {(0.0,0.002) (0.1,0.002) (0.2,0.002) (0.3,0.002) (0.4,0.002) (0.5,0.002) (0.6,0.002) (0.7,0.002) (0.8,0.002) (0.9,0.002) (1.0,0.002) (1.1,0.083) (1.2,0.165) (1.3,0.242) (1.4,0.313) (1.5,0.378) (1.6,0.437) (1.7,0.491) (1.8,0.539) (1.9,0.583) (2.0,0.623) (2.1,0.643) (2.2,0.665) (2.3,0.688) (2.4,0.71) (2.5,0.732) (2.6,0.753) (2.7,0.773) (2.8,0.792) (2.9,0.81) (3.0,0.826) (3.2,0.855) (3.4,0.88) (3.6,0.9) (3.8,0.917) (4.0,0.931) (4.4,0.953) (4.8,0.967) (5.2,0.977) (5.6,0.984) (6.0,0.988) (6.5,0.992) (7.0,0.994)};
\draw[very thick,gray!60!black,densely dashed] plot coordinates {(0.0,0.002) (0.1,0.002) (0.2,0.002) (0.3,0.002) (0.4,0.002) (0.5,0.002) (0.6,0.002) (0.7,0.002) (0.8,0.002) (0.9,0.002) (1.0,0.002) (1.1,0.083) (1.2,0.165) (1.3,0.242) (1.4,0.313) (1.5,0.378) (1.6,0.358) (1.7,0.336) (1.8,0.313) (1.9,0.291) (2.0,0.269) (2.2,0.228) (2.4,0.191) (2.6,0.159) (2.8,0.133) (3.0,0.11) (3.2,0.091) (3.4,0.076) (3.6,0.063) (3.8,0.052) (4.0,0.043) (4.4,0.03) (4.8,0.021) (5.2,0.015) (5.6,0.011) (6.0,0.008) (6.5,0.006) (7.0,0.004)};
\node[font=\scriptsize,blue!60!black,anchor=west] at (3.3,0.8) {приток $1\tau$: записано};
\node[font=\scriptsize,gray!50!black,anchor=west] at (2.3,0.3) {приток $0.5\tau$: погасло};
\end{scope}
\end{tikzpicture}
\caption{Группа, сильно связанная сама с собой. (а)~Много \nt{ГЛУТ}-нейронов, возбуждающих друг друга, ведут себя как один нейрон с частотой $r$, который возбуждает сам себя с силой $w$; извне приходит приток $I$. (б)~Частота группы во времени при $w=1$ и той же кривой $f$, что на рис.~\ref{fig:bistable}. Приток $I=0.6$ подаётся на время, отмеченное полоской под осью. Если он держится $1\tau$, группа успевает перейти неустойчивую точку $r=0.5$, разгорается и остаётся гореть, когда приток кончился. Если только $0.5\tau$, она не дотягивает до этой точки и гаснет обратно: чтобы записать бит, приток должен длиться дольше примерно $0.7\tau$.}
\label{fig:recurrentgroup}
\end{figure}

\begin{figure}[t]
\centering
\begin{tikzpicture}[>=Stealth,x=5cm,y=3.2cm]
\draw[->,gray!70!black] (0,0) -- (1.1,0) node[below left,black] {\small $r$};
\draw[->,gray!70!black] (0,0) -- (0,1.1);
\draw[thick,gray!60!black] (0,0) -- (1.05,1.05) node[right,black] {\small $r$};
\draw[very thick,blue!60!black] plot[domain=0:1.05,samples=80] (\x,{1/(1+exp(-(\x-0.5)/0.08))});
\draw[very thick,red!70!black,densely dashed] plot[domain=0:1.05,samples=80] (\x,{1/(1+exp(-(\x+0.3-0.5)/0.08))});
\fill[blue!60!black] (0.002,0.002) circle (0.018);
\draw[blue!60!black,fill=white,thick] (0.5,0.5) circle (0.018);
\fill[blue!60!black] (0.998,0.998) circle (0.018);
\node[below right,font=\small] at (0.02,0.0) {выкл.};
\node[right,font=\small] at (0.52,0.46) {неустойчиво};
\node[below,font=\small] at (0.99,0.96) {вкл.};
\node[anchor=east,font=\small,red!70!black] at (0.19,0.62) {$f(wr+I)$};
\node[anchor=west,font=\small,blue!60!black] at (0.45,0.1) {$f(wr)$};
\end{tikzpicture}
\caption{Память из группы \nt{ГЛУТ}-нейронов с сильной обратной связью. Сплошная кривая --- без внешнего притока: два устойчивых пересечения с прямой $r$, «выключено» и «включено», и неустойчивое между ними. Короткий приток $I$ (штриховая кривая) оставляет только верхнее.}
\label{fig:bistable}
\end{figure}

Если обратная связь сильная, пересечений три: нижнее --- тишина, верхнее --- группа горит на полную, среднее неустойчиво. Получился элемент с двумя устойчивыми состояниями, триггер. Записать в него единицу легко: короткий приток поднимает кривую, нижнее пересечение исчезает, и группа разгорается --- и остаётся гореть, когда приток кончился (рис.~\ref{fig:recurrentgroup}б). Слишком короткий приток не успевает довести группу до неустойчивой точки, и она гаснет обратно. Такая самоподдерживающаяся активность, которая держится секунды после того, как раздражитель исчез, действительно наблюдается в мозге и считается основой кратковременной памяти~\cite{Wang2001}. Но это не единственный способ помнить секунды. Запись можно хранить и в медленной переменной самих синапсов: после пачки импульсов облегчение, как у синапса-«тире» главы~\ref{sec:code}, держится около секунды, и сеть между короткими вспышками может почти молчать --- не триггер, а заряженный конденсатор~\cite{Mongillo2008}. Такие «тихие» промежутки при удержании в памяти наблюдаются~\cite{Stokes2015}, и хранение без постоянных импульсов дешевле по энергии. Какой из двух способов кора использует в каком случае, обсуждается.

А как стереть? По утверждению~\ref{prop:monotone} активностью --- никак; остаётся что-то \emph{убрать}. Первый способ --- убрать поддержку: если группа горит только вместе с притоком от другой группы, память гаснет вместе с ним. Второй --- усталость. У настоящих синапсов при долгой работе запас медиатора истощается~\cite{Tsodyks1997}, а у нейронов нарастают медленные токи, снижающие возбудимость. Обратная связь слабеет, верхнее пересечение исчезает, и группа гаснет сама через секунды. Получается память, которая включается сигналом, а выключается только временем.

\section{Последовательности}

Вместо часов можно иметь \emph{таймер, запускаемый событием}. Соединим группы \nt{ГЛУТ}-нейронов в цепочку: первая возбуждает вторую, вторая --- третью, и так далее. Внешний сигнал зажигает первую группу, и волна активности пробегает по цепочке, каждое звено через одну-две задержки после предыдущего и только один раз: сразу после импульса нейроны невосприимчивы, а волна уже ушла дальше.

Такая цепочка отмеряет время от события: сработало звено номер $k$ --- значит, с момента запуска прошло $k$ задержек. Её выходы можно подать на мышцы и получить последовательность движений, которая разворачивается сама. У певчих птиц во время пения отдельные \nt{ГЛУТ}-нейроны одного из отделов мозга срабатывают строго по очереди, каждый в свой момент песни, --- очень похоже на такую цепочку~\cite{Hahnloser2002}.

В отличие от часов, цепочка молчит, пока её не запустили, отрабатывает один раз и отмеряет время только для тех, кто к ней подключён. Общего такта у сети по-прежнему нет.

\section{Чего не хватает}

Из одних \nt{ГЛУТ}-нейронов, без торможения, получается многое. Но трёх вещей не хватает, и все три из-за утверждения~\ref{prop:monotone}.

\emph{Выбор.} Сеть должна выбирать одно направление, одно действие. Если два раздражителя почти равны, в сети рис.~\ref{fig:itd} сработают оба выхода, и подавить слабый в пользу сильного нечем.

\emph{Сброс.} Стереть память сигналом нельзя.

\emph{Устойчивость.} В сети, где связи возникают не по плану инженера, петли положительной обратной связи неизбежны, и активность в них может нарастать лавиной (глава~\ref{sec:neurontypes}). Единственный тормоз --- та же усталость, медленная и грубая.

Все три беды лечатся одним лекарством --- нейроном, который гасит импульс импульсом. Это \nt{ГАМК}, и нужен он не для того, чтобы вычислять, а для того, чтобы выбирать, сбрасывать и держать вычисление под контролем.

\section{Задачи}

\begin{exercise}[\lvlA]
\label{ex:02:1}
\emph{Ряд детекторов для звука.} Сигналы от ушей бегут по проводам рис.~\ref{fig:itd} со скоростью $v=5\,\text{м/с}$; наибольшая разница времён прихода --- $0.22\,\text{м}/343\,\text{м/с}$.
\par\noindent(а)~Какой длины $L$ должен быть ряд, чтобы покрыть все направления? С каким шагом ставить детекторы, чтобы соседние отвечали разницам $\delta t$, отличающимся на $10$~мкс, и сколько их нужно?
\par\noindent(б)~Детекторы --- нейроны~\eqref{eq:glutneuron} с $\tau=0.5\ms$ и $\theta=1.5w$, на каждый вход приходит по одному импульсу. Сколько детекторов сработает на один звук? Что это значит для утверждения «номер сработавшего детектора и есть направление» и как получатель может извлечь из ряда точность в $10$~мкс?
\end{exercise}

\begin{exercise}[\lvlB]
\label{ex:02:2}
\emph{Неравные входы смещают окно.} Детектор совпадений~\eqref{eq:glutneuron} получает по одному импульсу от двух входов с весами $w_1$ и $w_2$, $w_1,w_2<\theta<w_1+w_2$; второй импульс приходит на $\delta=t_2-t_1$ позже первого ($\delta$ может быть отрицательным).
\par\noindent(а)~Выведите окно совпадения: множество $\delta$, при которых нейрон срабатывает. Проверьте, что при $w_1=w_2$ получается~\eqref{eq:window}.
\par\noindent(б)~Найдите центр окна $c$ и вычислите окно и $c$ при $w_1=1$, $w_2=0.8$, $\theta=1.5$, $\tau=0.5\ms$. Какой вход должен прийти раньше, чтобы ответ был наиболее надёжным, и почему?
\par\noindent(в)~Во всём ряду рис.~\ref{fig:itd} вход от левого уха на $20\,\%$ слабее правого. Как это сказывается на найденном направлении, и сравните ошибку с разрешением около десяти микросекунд. При каком $\theta$ смещение исчезает и чем за это приходится платить?
\end{exercise}

\begin{exercise}[\lvlB]
\label{ex:02:3}
\emph{Монотонность в непрерывном времени.} Частоты подчиняются уравнениям $\tau\,\dd r_i/\dd t=-r_i+f_i\bigl(\sum_jw_{ij}r_j+I_i(t)\bigr)$, где $f_i$ не убывают и удовлетворяют условию Липшица с константой $\ell$.
\par\noindent(а)~Пусть $w_{ij}\ge0$ при $i\ne j$ (знак $w_{ii}$ любой). Докажите: если $r(0)\le r'(0)$ и $I(t)\le I'(t)$ покомпонентно, то $r(t)\le r'(t)$ при всех $t>0$. Указание: рассмотрите $z=r'-r$ и $\sum_i\max(-z_i,0)$.
\par\noindent(б)~Выведите отсюда утверждение~\ref{prop:monotone} для сети, стартующей из тишины, без итераций по шагам.
\par\noindent(в)~Знаки весов произвольны. Докажите, что замена $r_i\to s_ir_i$, $s_i=\pm1$, переводит систему в систему из пункта~(а) тогда и только тогда, когда в каждом цикле графа связей (без учёта направления стрелок, без петель $i\to i$) чётное число отрицательных связей. Примените это к двум группам с взаимным торможением и к петле «\nt{ГЛУТ} возбуждает \nt{ГАМК}, \nt{ГАМК} тормозит \nt{ГЛУТ}» из главы~\ref{sec:gaba}. Что говорит результат о том, какая из них может колебаться?
\end{exercise}

\begin{exercise}[\lvlB]
\label{ex:02:4}
\emph{Проектирование: двухпроводный сумматор.} Полный сумматор складывает три бита $x,y,z$: выдаёт бит суммы $S=x\oplus y\oplus z$ и бит переноса $C$ --- единицу, если единиц не меньше двух.
\par\noindent(а)~Докажите: нейрон с неотрицательными весами, получающий пары проводов $(x_i,\bar x_i)$, на полных входах вычисляет любую пороговую функцию $\bigl[\sum_ic_ix_i\ge\theta\bigr]$ с весами $c_i$ \emph{любого} знака.
\par\noindent(б)~Постройте двухпроводный полный сумматор (пары $(S,\bar S)$ и $(C,\bar C)$) не более чем из восьми \nt{ГЛУТ}-нейронов, указав веса и пороги. Указание: начните с нейронов $g_k=[x+y+z\ge k]$.
\par\noindent(в)~Входы приходят асинхронно: пока какая-то пара молчит, её значение ещё не пришло. Докажите, что в вашей схеме ни один провод ответа не сработает неверно ни при каком порядке прихода входов.
\par\noindent(г)~Сколько нейронов потребовала бы схема из элементов И и ИЛИ, как на рис.~\ref{fig:dualxor}?
\end{exercise}

\begin{exercise}[\lvlB]
\label{ex:02:5}
\emph{Моделирование: сколько длится запись.} Возьмите модель рис.~\ref{fig:recurrentgroup}: $\tau\,\dd r/\dd t=-r+f(wr+I(t))$, $f(u)=1/\bigl(1+e^{-(u-0.5)/0.08}\bigr)$, $w=1$; до записи группа в нижнем устойчивом состоянии $r_0$; приток $I$ подаётся на время $T$.
\par\noindent(а)~Интегрируя методом Эйлера (шаг $10^{-4}\tau$), воспроизведите рис.~\ref{fig:recurrentgroup}б и найдите $r_0$.
\par\noindent(б)~Найдите наименьшую длительность записи $T_{\min}(I)$ при $I$ от $0.25$ до $2$ и постройте её график. Чему равна $T_{\min}(0.6)$?
\par\noindent(в)~Докажите, что при любом $I$ $T_{\min}>\tau\ln\bigl[(1-r_0)/(1-r_u)\bigr]$, где $r_u=0.5$ --- неустойчивое равновесие, и что при $I\to\infty$ $T_{\min}$ стремится к этой границе.
\par\noindent(г)~Найдите $I_c$, ниже которого бит не записывается ни при какой длительности, и проверьте численно, как растёт $T_{\min}$ при $I\to I_c$. Какой смысл имеют оба предела для электронщика, привыкшего к триггерам?
\end{exercise}

\begin{exercise}[\lvlA]
\label{ex:02:6}
\emph{Цепочка как таймер.} Звено цепочки из раздела о последовательностях --- группа из $100$ \nt{ГЛУТ}-нейронов; задержка звена $2\ms$ со случайным разбросом $0.2\ms$, независимым от звена к звену.
\par\noindent(а)~Сколько нужно звеньев и нейронов, чтобы отмерить $1\,\text{с}$? Сколько импульсов тратится на один отсчёт?
\par\noindent(б)~Каков разброс момента, когда срабатывает последнее звено, и его относительная величина? То же для $100\ms$. Как относительная ошибка зависит от отмеряемого интервала?
\par\noindent(в)~Допустим, измерение показывает, что относительная ошибка таймера не зависит от интервала и равна $5\,\%$. Какой источник разброса в цепочке дал бы такой результат?
\end{exercise}

\begin{exercise}[\lvlC]
\label{ex:02:7}
\emph{Исследование: триггер или конденсатор.} В разделе о памяти остался открытым вопрос, как удерживать бит секунды: самоподдерживающейся активностью группы или медленной переменной синапсов, которую изредка освежают короткие пачки.
\par\noindent(а)~Группа из $N=100$ нейронов удерживает бит $10\,\text{с}$. Триггер: все нейроны работают с частотой $20$ импульсов в секунду. Конденсатор: облегчение $u$ спадает с постоянной времени $\tau_F=1\,\text{с}$, бит читается, пока $u\ge u_{\max}/2$, и каждое освежение --- пачка из трёх импульсов у каждого нейрона. Найдите частоту освежений, число импульсов и энергию при стоимости импульса $10^{-10}$~Дж (глава~\ref{sec:code}).
\par\noindent(б)~На $200\ms$ всю группу глушат общим сигналом. Что происходит с битом в каждой модели?
\par\noindent(в)~Открытая часть. Постройте обе модели в виде уравнений для частот (для второй добавьте уравнение для $u$, умножающего вес обратной связи) и предложите опыт с короткими неспецифическими импульсами на группу, который различает две модели по отклику, а не по фоновой активности. Какие измерения внутри одной и той же модели мешают сделать вывод, и какие результаты допускают смешанный механизм?
\end{exercise}
